\chapter{Controlled-Collision Construction Benchmarks}
\label{app:controlled_comparison}

After implementing boundary conditions, we ran into very slow simulations that were impractical to run.
Analyzing the runtime, we managed to isolate the bottleneck to the controlled unitaries in the collision step.
On a whim we decided to try implementing the controlled gate by manually building the matrix via Kronecker-product construction and apply it as a single \texttt{UnitaryGate}.
To our surprise, this sped up simulation quite a bit.

Performing a detailed benchmark revealed the explanation: \texttt{UnitaryGate} is considered elementary in the default Aer Simulator.
This means it won't get transpiled and it will be applied via simple vector-matrix multiplication.
Transpiling against a more realistic elementary gate set, like the IBM Eagle's, shows that the two approaches are almost identical in gate count and circuit depth.
Surprisingly, for diagonal gates, the manual approach still appeared to result in a 4x decrease in depth.

\begin{table}[H]
\centering
\caption{Gate count / depth transpilation comparison and runtime}
\label{tab:controlled_comparison}
\resizebox{\textwidth}{!}{%
\begin{tabular}{l c r r r r r r}
\toprule
 & & \multicolumn{2}{c}{Eagle basis} & \multicolumn{2}{c}{Default (Aer) basis} & \multicolumn{2}{c}{Runtime (ms)} \\
\cmidrule(lr){3-4} \cmidrule(lr){5-6} \cmidrule(lr){7-8}
Operator & Qubits & \texttt{.ctrl} & manual & \texttt{.ctrl} & manual & \texttt{.ctrl} & manual \\
\midrule
\multicolumn{8}{l}{\emph{Diagonal} (collision without BC)} \\
Random phase diagonal & 4 & 281 / 206       & 66 / 55       & 46 / 43        & 1 / 1 & 2.0    & 0.7     \\
Random phase diagonal & 6 & 1,197 / 916     & 258 / 241     & 190 / 185      & 1 / 1 & 4.3    & 1.8     \\
Random phase diagonal & 8 & 4,849 / 3,786   & 1,026 / 1,004 & 766 / 759      & 1 / 1 & 14.9   & 28.8    \\
\midrule
\multicolumn{8}{l}{\emph{General unitary} (collision combined with BC, post-SVD)} \\
Random unitary   & 4 & 1,158 / 724     & 1,156 / 724   & 635 / 467      & 1 / 1 & 11.9   & 0.6     \\
Random unitary   & 6 & 19,208 / 12,292 & 19,206 / 12,292 & 10,845 / 8,195 & 1 / 1 & 201.7  & 1.0     \\
Random unitary   & 8 & 310,282 / 199,684 & 310,279 / 199,683 & 176,479 / 134,147 & 1 / 1 & 3,666.1 & 19.0 \\
\bottomrule
\end{tabular}%
}
\end{table}
